\textbf{Phần 1. Trắc nghiệm lựa chọn}
\setcounterpageref{ex}{0}

\begin{ex} %[5H3N16]
Trong bảng đơn vị đo diện tích, mỗi đơn vị đo gấp bao nhiêu lần đơn vị bé hơn tiếp liền?
\choice
{$10$ lần}
{$1\,000$ lần}
{\True $100$ lần}
{$10\,000$ lần}
\loigiai{
Trong bảng đơn vị đo diện tích, mỗi đơn vị lớn gấp $100$ lần đơn vị bé hơn tiếp liền.\\
\textbf{Đáp án của câu hỏi là C}
}
\end{ex}

\begin{ex} %[5H3N16]
Đổi $4\text{ m}^2$ ra đề-xi-mét vuông ta được kết quả là
\choice
{\True $400\text{ dm}^2$}
{$40\text{ dm}^2$}
{$4\,000\text{ dm}^2$}
{$\dfrac{4}{100}\text{ dm}^2$}
\loigiai{
Vì $1\text{ m}^2 = 100\text{ dm}^2$ nên $4\text{ m}^2 = 400\text{ dm}^2$.\\
\textbf{Đáp án của câu hỏi là A}
}
\end{ex}

\begin{ex} %[5H3H16]
Đổi $2\text{ m}^2\ 5\text{ dm}^2$ ra đề-xi-mét vuông ta được kết quả là
\choice
{$25\text{ dm}^2$}
{\True $205\text{ dm}^2$}
{$250\text{ dm}^2$}
{$2\,005\text{ dm}^2$}
\loigiai{
Ta có: $2\text{ m}^2\ 5\text{ dm}^2 = 2 \times 100 + 5 = 205\text{ dm}^2$.\\
\textbf{Đáp án của câu hỏi là B}
}
\end{ex}

\begin{ex} %[5H3H16]
Đổi $40\text{ cm}^2\ 4\text{ mm}^2$ ra mi-li-mét vuông ta được kết quả là
\choice
{$44\text{ mm}^2$}
{$404\text{ mm}^2$}
{\True $4\,004\text{ mm}^2$}
{$40\,004\text{ mm}^2$}
\loigiai{
Ta có: $40\text{ cm}^2\ 4\text{ mm}^2 = 40 \times 100 + 4 = 4\,004\text{ mm}^2$.\\
\textbf{Đáp án của câu hỏi là C}
}
\end{ex}

\begin{ex} %[5H3H16]
Viết số đo $3\text{ m}^2\ 7\text{ dm}^2$ dưới dạng số thập phân có đơn vị mét vuông ta được
\choice
{$3{,}7\text{ m}^2$}
{$3{,}007\text{ m}^2$}
{$37{,}0\text{ m}^2$}
{\True $3{,}07\text{ m}^2$}
\loigiai{
Ta có: $3\text{ m}^2\ 7\text{ dm}^2 = 3\dfrac{7}{100}\text{ m}^2 = 3{,}07\text{ m}^2$.\\
\textbf{Đáp án của câu hỏi là D}
}
\end{ex}

\begin{ex} %[5H3H16]
Viết số đo $12\text{ km}^2\ 50\text{ ha}$ dưới dạng số thập phân có đơn vị ki-lô-mét vuông ta được
\choice
{\True $12{,}5\text{ km}^2$}
{$12{,}05\text{ km}^2$}
{$12{,}005\text{ km}^2$}
{$125{,}0\text{ km}^2$}
\loigiai{
Ta có: $12\text{ km}^2\ 50\text{ ha} = 12\dfrac{50}{100}\text{ km}^2 = 12{,}50\text{ km}^2 = 12{,}5\text{ km}^2$.\\
\textbf{Đáp án của câu hỏi là A}
}
\end{ex}

\begin{ex} %[5H3H16]
Số thích hợp điền vào chỗ chấm để $615\text{ dm}^2 = \dots\text{ m}^2\ \dots\text{ dm}^2$ là
\choice
{$6\text{ m}^2\ 5\text{ dm}^2$}
{\True $6\text{ m}^2\ 15\text{ dm}^2$}
{$61\text{ m}^2\ 5\text{ dm}^2$}
{$60\text{ m}^2\ 15\text{ dm}^2$}
\loigiai{
Ta có: $615\text{ dm}^2 = 600\text{ dm}^2 + 15\text{ dm}^2 = 6\text{ m}^2\ 15\text{ dm}^2$.\\
\textbf{Đáp án của câu hỏi là B}
}
\end{ex}

\begin{ex} %[5H3H16]
Số đo thích hợp nhất chỉ diện tích mặt bàn học sinh là
\choice
{$50\text{ cm}^2$}
{$50\text{ m}^2$}
{$50\text{ ha}$}
{\True $50\text{ dm}^2$}
\loigiai{
Mặt bàn học sinh có diện tích thích hợp nhất là $50\text{ dm}^2$ (tương đương $0{,}5\text{ m}^2$).\\
\textbf{Đáp án của câu hỏi là D}
}
\end{ex}

\begin{ex} %[5H3H16]
Số đo diện tích thích hợp nhất cho mặt sàn một căn hộ chung cư là
\choice
{$225\text{ cm}^2$}
{\True $94\text{ m}^2$}
{$94\text{ dm}^2$}
{$8\,310\text{ km}^2$}
\loigiai{
Diện tích mặt sàn căn hộ thường đo bằng mét vuông, phù hợp nhất là $94\text{ m}^2$.\\
\textbf{Đáp án của câu hỏi là B}
}
\end{ex}

\begin{ex} %[5H3V16]
Ô cửa sổ thông gió ở bếp nhà Mai có dạng hình chữ nhật với chiều dài $1\text{ m}$ và chiều rộng $\dfrac{3}{5}\text{ m}$. Diện tích của ô cửa sổ đó tính theo đề-xi-mét vuông là
\choice
{$\dfrac{3}{5}\text{ dm}^2$}
{$6\text{ dm}^2$}
{$600\text{ dm}^2$}
{\True $60\text{ dm}^2$}
\loigiai{
Đổi: $1\text{ m} = 10\text{ dm}$; $\dfrac{3}{5}\text{ m} = 6\text{ dm}$.\\
Diện tích ô cửa sổ là: $10 \times 6 = 60\text{ (dm}^2\text{)}$.\\
\textbf{Đáp án của câu hỏi là D}
}
\end{ex}

\begin{ex} %[5H3V16]
Bác Tư có $1\text{ ha}$ đất. Bác đã dùng $\dfrac{9}{10}$ diện tích mảnh đất để trồng cây và làm đường, phần còn lại để xây nhà. Diện tích đất dùng để xây nhà là
\choice
{\True $1\,000\text{ m}^2$}
{$100\text{ m}^2$}
{$9\,000\text{ m}^2$}
{$10\,000\text{ m}^2$}
\loigiai{
Đổi $1\text{ ha} = 10\,000\text{ m}^2$.\\
Phần diện tích đất xây nhà chiếm: $1 - \dfrac{9}{10} = \dfrac{1}{10}$ (diện tích khu đất).\\
Diện tích đất để xây nhà là: $10\,000 \times \dfrac{1}{10} = 1\,000\text{ (m}^2\text{)}$.\\
\textbf{Đáp án của câu hỏi là A}
}
\end{ex}

\begin{ex} %[5H3V16]
Một mảnh vườn có diện tích $2\,500\text{ dm}^2$. Người ta dùng $\dfrac{2}{5}$ diện tích vườn để trồng hoa. Diện tích trồng hoa theo đơn vị mét vuông là
\choice
{$100\text{ m}^2$}
{$1\,000\text{ m}^2$}
{\True $10\text{ m}^2$}
{$15\text{ m}^2$}
\loigiai{
Đổi $2\,500\text{ dm}^2 = 25\text{ m}^2$.\\
Diện tích trồng hoa là: $25 \times \dfrac{2}{5} = 10\text{ (m}^2\text{)}$.\\
\textbf{Đáp án của câu hỏi là C}
}
\end{ex}

\textbf{Phần 2. Trắc nghiệm đúng sai}
\setcounterpageref{ex}{0}

\begin{ex} %[5H3H16]
Xét tính đúng sai của các khẳng định sau về đổi đơn vị đo diện tích:
\choiceTF
{\True $12\text{ m}^2 = 1\,200\text{ dm}^2$.}
{$2\,500\text{ dm}^2 = 250\text{ m}^2$.}
{\True $8\text{ cm}^2\ 10\text{ mm}^2 = 8{,}1\text{ cm}^2$.}
{$30\text{ dm}^2 = 3\text{ m}^2$.}
\loigiai{
\begin{itemize}
    \item $12\text{ m}^2 = 12 \times 100 = 1\,200\text{ dm}^2$.
    \item $2\,500\text{ dm}^2 = 2\,500 : 100 = 25\text{ m}^2 \neq 250\text{ m}^2$.
    \item $8\text{ cm}^2\ 10\text{ mm}^2 = 8\dfrac{10}{100}\text{ cm}^2 = 8{,}1\text{ cm}^2$.
    \item $30\text{ dm}^2 = \dfrac{30}{100}\text{ m}^2 = 0{,}3\text{ m}^2 \neq 3\text{ m}^2$.
\end{itemize}
\begin{itemize}
    \item a) \textbf{Đúng}.
    \item b) \textbf{Sai}: $2\,500\text{ dm}^2 = 25\text{ m}^2$.
    \item c) \textbf{Đúng}.
    \item d) \textbf{Sai}: $30\text{ dm}^2 = 0{,}3\text{ m}^2$.
\end{itemize}
}
\end{ex}

\begin{ex} %[5H3V16]
Bác Tư có thửa ruộng diện tích $1\text{ ha}$. Bác dành $\dfrac{9}{10}$ diện tích để trồng lúa và làm mương nước, phần còn lại xây trang trại. Xét tính đúng sai của các khẳng định sau:
\choiceTF
{\True Diện tích cả thửa ruộng là $10\,000\text{ m}^2$.}
{Diện tích đất trồng lúa và làm mương nước là $900\text{ m}^2$.}
{\True Diện tích xây trang trại là $1\,000\text{ m}^2$.}
{\True Diện tích đất trồng lúa và làm mương nước gấp $9$ lần diện tích đất xây trang trại.}
\loigiai{
\begin{itemize}
    \item $1\text{ ha} = 10\,000\text{ m}^2$.
    \item Diện tích trồng lúa và mương nước: $10\,000 \times \dfrac{9}{10} = 9\,000\text{ m}^2 \neq 900\text{ m}^2$.
    \item Diện tích xây trang trại: $10\,000 - 9\,000 = 1\,000\text{ m}^2$.
    \item Tỉ số: $9\,000 : 1\,000 = 9$ lần.
\end{itemize}
\begin{itemize}
    \item a) \textbf{Đúng}.
    \item b) \textbf{Sai}: Diện tích là $9\,000\text{ m}^2$.
    \item c) \textbf{Đúng}.
    \item d) \textbf{Đúng}.
\end{itemize}
}
\end{ex}

\textbf{Phần 3. Trả lời ngắn}
\setcounterpageref{ex}{0}

\begin{ex} %[5H3H16]
Đổi $2\,500\text{ dm}^2$ ra mét vuông.
\shortans{25}
\loigiai{
\textbf{Đáp số:} $25$\\
$2\,500\text{ dm}^2 = 2\,500 : 100 = 25\text{ m}^2$.
}
\end{ex}

\begin{ex} %[5H3H16]
Đổi $40\text{ cm}^2\ 4\text{ mm}^2$ ra mi-li-mét vuông.
\shortans{4004}
\loigiai{
\textbf{Đáp số:} $4004$\\
$40\text{ cm}^2\ 4\text{ mm}^2 = 40 \times 100 + 4 = 4\,004\text{ mm}^2$.
}
\end{ex}

\begin{ex} %[5H3H16]
Viết số đo $30\text{ dm}^2$ dưới dạng số thập phân có đơn vị mét vuông. (Nhập kết quả dưới dạng số thập phân).
\shortans{0,3}
\loigiai{
\textbf{Đáp số:} $0{,}3$\\
$30\text{ dm}^2 = \dfrac{30}{100}\text{ m}^2 = 0{,}3\text{ m}^2$.
}
\end{ex}

\begin{ex} %[5H3H16]
Viết số đo $12\text{ km}^2\ 50\text{ ha}$ dưới dạng số thập phân có đơn vị ki-lô-mét vuông. (Nhập kết quả dưới dạng số thập phân).
\shortans{12,5}
\loigiai{
\textbf{Đáp số:} $12{,}5$\\
$12\text{ km}^2\ 50\text{ ha} = 12{,}5\text{ km}^2$.
}
\end{ex}

\begin{ex} %[5H3V16]
Ô cửa sổ hình chữ nhật có chiều dài $1\text{ m}$, chiều rộng $\dfrac{3}{5}\text{ m}$. Hỏi diện tích ô cửa sổ đó là bao nhiêu đề-xi-mét vuông?
\shortans{60}
\loigiai{
\textbf{Đáp số:} $60$\\
$1\text{ m} = 10\text{ dm}$; $\dfrac{3}{5}\text{ m} = 6\text{ dm} \Rightarrow S = 10 \times 6 = 60\text{ dm}^2$.
}
\end{ex}

\begin{ex} %[5H3V16]
Bác Tư có thửa đất $1\text{ ha}$. Bác dùng $\dfrac{9}{10}$ thửa đất làm vườn, phần còn lại xây nhà. Hỏi diện tích xây nhà của bác Tư là bao nhiêu mét vuông?
\shortans{1000}
\loigiai{
\textbf{Đáp số:} $1000$\\
Diện tích đất xây nhà là: $10\,000 \times \left(1 - \dfrac{9}{10}\right) = 1\,000\text{ (m}^2\text{)}$.
}
\end{ex}
